\section{A boundary from finite-energy functions}
\label{sec:compactification}

We now construct the compact space used in the boundary argument.
The inputs are the group $\Gamma$, its triangular Cayley $2$-complex
$\mathcal C$, its graph $G$, and the proper closed inclusion
$\mathcal D_0\subsetneq\mathcal D$ obtained above.  Thus $G$ is connected,
locally finite, and one-ended; $\mathcal C$ is simply connected; and the
number of cells incident to an edge is uniformly bounded.  We retain
cellular edges, including parallel edges, in all energy sums.
The output will be a compact metrizable space $X=\Gamma\sqcup B$, with a
nondegenerate connected boundary $B$.  In addition, vertices approaching
the same boundary point will admit joining paths uniformly near that
point.  This last property concerns vertex paths in $G$; local
connectivity of $B$ is a separate question.

\subsection{A function with connected level sides}

For a real function $F$ on the vertices, write
\[
 \mathcal E(F)=\sum_e |dF(e)|^2,
 \qquad dF(e)=F(e^+)-F(e^-),
\]
where one orientation of each cellular edge is used.  We call $F$ a
finite-energy function when $\mathcal E(F)<\infty$.  A level $t$ is
\emph{regular for $F$} if it is not a vertex value.  At such a level put
\[
 S(F,t)=\{e:\min(F(e^-),F(e^+))<t<
                       \max(F(e^-),F(e^+))\}.
\]
These are the edges crossing the level.  Two edges in a set of crossing
edges are \emph{dually adjacent} if they occur on the boundary of a common
$2$-cell.  Dual components are components for this adjacency relation.

\begin{lemma}\label{compact:potential}
There exists a finite-energy function $f:\Gamma\to[0,1]$ which does not
tend to a single constant outside finite subsets of $\Gamma$, and for
which both $\{f<t\}$ and $\{f>t\}$ are connected vertex subgraphs whenever
$0<t<1$ is regular for $f$.  Every such $S(f,t)$ is dually connected.
\end{lemma}

\begin{proof}
Fix a vertex $a$.  For each $b\ne a$, the functional
\[
 L_b(df)=f(b)-f(a)
\]
is well defined on $\mathcal D$.  A finite path from $a$ to $b$ of length
$\ell$ gives $|L_b(df)|\le\sqrt{\ell}\,\|df\|_2$.  The functional is
nonzero, as one sees by testing the gradient of the singleton indicator
of $b$.  Let $r_b\in\mathcal D$ be its Riesz representer.  The unique
gradient of minimum norm subject to $L_b(df)=1$ is
\[
 df_b=\frac{r_b}{\|r_b\|_2^2};
\]
normalize its potential by $f_b(a)=0$.

The vectors $r_b$ span densely in $\mathcal D$.  Indeed, a gradient
orthogonal to all of them has a potential constant at every vertex and
is therefore zero.  Since $\mathcal D_0$ is a proper closed subspace,
some $df_b$ lies outside $\mathcal D_0$.  Fix this minimizer and write it
as $df$.  Replacing $f$ by $\min(1,\max(0,f))$ preserves the values $0,1$
at $a,b$ and cannot increase any edge energy.  Uniqueness of the
minimizer implies $0\le f\le1$.

Suppose $f$ tends to a constant $c$ outside finite vertex sets.  For
$\epsilon>0$ define
\[
 u_\epsilon(v)=\operatorname{sgn}(f(v)-c)
                      \max(|f(v)-c|-\epsilon,0).
\]
Each $u_\epsilon$ has finite support.  The function
$(f-c)-u_\epsilon$ is the truncation of $f-c$ to
$[-\epsilon,\epsilon]$, so its edge differences are bounded in absolute
value by $|df|$ and tend pointwise to zero.  Dominated convergence gives
$du_\epsilon\to df$ in edge $\ell^2$, contrary to
$df\notin\mathcal D_0$.

Fix a regular $t\in(0,1)$.  If a component $U$ of $\{f>t\}$ misses $b$,
replace $f$ by $t$ on $U$.  Both prescribed values are preserved.
Internal edge energies decrease to zero.  At a boundary edge with one
endpoint in $U$, the other endpoint has value below $t$, so its energy
decreases strictly.  Such an edge exists because $G$ is connected and
$a\notin U$.  This modification is legitimate even when $U$ is
infinite: each edge energy weakly decreases, and at least one decreases
by a positive amount.  It contradicts minimality.  Hence every upper
component contains $b$, and the upper side is connected.  The lower
side is connected by the same argument with $a$.

Let $C$ be a dual component of $S(f,t)$, and restrict the cochain
$d1_{\{f>t\}}$ to $C$, setting it to zero on other edges.  In any cell,
all crossing edges lie in one dual component.  The signed sum of the
restricted cochain around that cell is therefore either the original
zero sum or an entirely zero sum.  This remains true when an edge has
several occurrences on the cell boundary.  Simple connectivity of
$\mathcal C$ gives a potential for the restricted cochain.  That
potential is constant on each of the two connected level sides, and
an edge of $C$ shows that the constants differ.  It consequently has
nonzero difference across every edge of $S(f,t)$, so no crossing edge
could have been omitted.  Thus $C=S(f,t)$.
\end{proof}

For a function $F$ and $g\in\Gamma$, define its translate by
$(gF)(v)=F(g^{-1}v)$.  Let $\mathcal F_0$ consist of all translates of
$f$, together with the singleton indicators $1_{\{v\}}$ for
$v\in\Gamma$.  These are countably many $[0,1]$-valued finite-energy
functions.

\begin{lemma}\label{compact:wall-components}
Every regular cut of a coordinate in $\mathcal F_0$ has finitely many
dual components.  A nonempty regular cut of a translate of $f$ is
dually connected, and a regular cut of a singleton indicator is finite.
\end{lemma}

\begin{proof}
The assertion for translates follows from Lemma~\ref{compact:potential};
levels outside $[0,1]$ have empty cuts.  The only nonempty cuts of a
singleton indicator consist of edges incident to its vertex, a finite
set by local finiteness.
\end{proof}

\subsection{Boxes with finitely many components}

A box in the base-coordinate cube $[0,1]^{\mathcal F_0}$ is specified by
finitely many conditions
\[
 O=\{(z_F):a_i<z_{F_i}<b_i,\quad 1\le i\le m\},
 \qquad F_i\in\mathcal F_0.
\]
Endpoints outside $[0,1]$ are allowed.  We retain the defining
inequalities as part of the box data and translate them without
changing their endpoints.  We use the same symbol for its vertex set,
obtained by evaluating the coordinates.  A box is
\emph{good} if all its endpoints are regular and any two of its endpoint
cuts have only finitely many pairs of edges which coincide or occur on
a common cell.

\begin{lemma}\label{compact:good-boxes}
Good boxes form a basis of the base-coordinate cube.  Every good box
has finitely many graph components on its vertex set.  There is a
countable $\Gamma$-invariant basis of good boxes.
\end{lemma}

\begin{proof}
Write $e\sim e'$ when $e=e'$ or the edges occur on a common cell.  This
symmetric relation has degree bounded by a constant $D$, by bounded
cell incidence and bounded cell length.  For finite-energy coordinates
$F,F'$, Tonelli's theorem and Cauchy--Schwarz give
\begin{align*}
 &\int_{\mathbb R^2}
 \#\{(e,e'):e\sim e',\ e\in S(F,s),\ e'\in S(F',t)\}\,ds\,dt\\
 &\hspace{15mm}=\sum_{e\sim e'}|dF(e)|\,|dF'(e')|
 \le D\|dF\|_2\|dF'\|_2<\infty.
\end{align*}
Thus almost every pair of endpoint parameters has finitely many
interactions.  The argument also applies when $F=F'$, since the two
parameters remain independent.  For a fixed finite list of box
conditions, intersect the finitely many full-measure conditions for
endpoint pairs and exclude the countably many vertex values.  The
resulting endpoint tuples are dense in their parameter domain.  They
therefore supply arbitrarily small good boxes about every point of
the cube.

Fix a good box and list its endpoint cuts as $S_1,\ldots,S_N$.
Delete all of their edges from $G$, obtaining a graph $G'$.  On every
component of $G'$ all endpoint signs are constant.  Let $E$ be the
finite set of cut edges which also occur in another cut or share a
cell with an edge of another cut.  For each $i$, the dual graph on
$S_i\setminus E$ still has finitely many components.  Indeed, its
original dual graph has finitely many components by
Lemma~\ref{compact:wall-components}, has bounded degree, and has had
only finitely many vertices deleted.

Consider two dually adjacent edges of $S_i\setminus E$.  On their
common triangular cell, their endpoints on a specified side of cut
$i$ either coincide or are joined by the remaining cell edge.  This
joining edge is not in $S_i$.  Moreover, the cell contains no edge of
any other endpoint cut, by the definition of $E$.  The joining edge
therefore lies in $G'$.  Following a finite dual path shows that all
endpoints on either specified side of a dual component of
$S_i\setminus E$ belong to one component of $G'$.  Such a dual
component is incident to at most two components of $G'$.  Each edge
of $E$ is likewise incident to at most two.

If any cut is nonempty, every component of $G'$ is incident to a cut
edge, since the original graph is connected.  The preceding finite
count proves that $G'$ has finitely many components.  If all cuts are
empty, $G'=G$ is connected.  The vertex set of $O$ is a union of those
components with the specified signs; its induced subgraph has exactly
these components.  This proves the required finiteness.  The argument
also covers repeated vertices or edges on a triangular boundary: the
same-side joining path then either has length zero or uses the
remaining edge.

There are countably many finite lists of base coordinates.  For each
list, choose a countable dense set of good endpoint tuples.  The
resulting boxes form a countable basis.  Add all their translates;
this remains countable, and goodness is preserved by the action.
\end{proof}

Fix such an invariant countable basis, denoted $\mathcal O$.  A box can
have several components even though each individual level side is
connected.  We next add coordinates which remember the component of
each box.  The preceding finiteness will make that information
persist at boundary points.

\subsection{Coordinates which distinguish box components}

For $O\in\mathcal O$, with conditions $a_i<F_i<b_i$, put
\[
 \psi_O=\max\bigl(0,\min\bigl(\{1\}\cup
                   \{F_i-a_i,b_i-F_i:1\le i\le m\}\bigr)\bigr).
\]
This is a continuous function of the base coordinates, with values in
$[0,1]$, positive precisely inside $O$.  For each graph component $U$
of the vertex set of $O$, define on vertices
\begin{equation}\label{compact:localized}
 h_{O,U}=\psi_O1_U.
\end{equation}
Let $\mathcal F$ be the family consisting of $\mathcal F_0$ and all
these localized coordinates.  It is countable and invariant under
$\Gamma$, since translation carries components of a box to components
of its translated box.

\begin{lemma}\label{compact:localized-energy}
Every member of $\mathcal F$ has finite energy.  For almost every
$t\in(0,1)$,
\begin{equation}\label{compact:shifted-walls}
 S(h_{O,U},t)\subseteq
 \bigcup_{i=1}^m
       \bigl(S(F_i,a_i+t)\cup S(F_i,b_i-t)\bigr),
\end{equation}
where every displayed level is regular for its coordinate.
\end{lemma}

\begin{proof}
The operations defining the margin give
\[
 |d\psi_O(e)|\le\max_{1\le i\le m}|dF_i(e)|.
\]
Two adjacent vertices with positive margin belong to the same
component of $O$.  Thus either both endpoints contribute their margin
to $h_{O,U}$, both contribute zero, or the endpoint outside $U$ has
zero margin.  In every case
\[
 |dh_{O,U}(e)|\le |d\psi_O(e)|,
 \qquad \|dh_{O,U}\|_2^2\le\sum_{i=1}^m\|dF_i\|_2^2.
\]
If there are no conditions, the margin and localized coordinate are
constant and their gradients vanish.

For the cut inclusion, the endpoint where $h_{O,U}>t$ belongs to $U$
and satisfies every margin inequality strictly above $t$.  At the
other endpoint, either it still belongs to $U$ and some margin is
below $t$, or it lies outside $U$.  In the latter case its margin must
be zero by adjacency, and again one of the margin inequalities fails.
Since $t<1$, the constant cap in the definition of $\psi_O$ contributes
no additional cut.  The crossed margin is one of the levels
$a_i+t$ or $b_i-t$, proving \eqref{compact:shifted-walls}.  All
exceptional values, for which one of these levels or $t$ itself is a
vertex value of its coordinate, form a countable set.
\end{proof}

\subsection{The compact space and its boundary}

Compactifying a graph by bounded finite-energy functions is the Royden
method \cite[Section~4]{KellerLenzSchmidtWojciechowski2017}. We use a
selected countable family, with localized coordinates that record
components of boxes.

Map each $v\in\Gamma$ to its profile $(F(v))_{F\in\mathcal F}$ and
let $X$ be the closure of these profiles in $[0,1]^{\mathcal F}$.  The
singleton indicators make this map injective, so we identify vertices
with their profiles.  They also make each vertex isolated in $X$:
the coordinate $1_{\{v\}}$ equals $1$ only at $v$, even after taking
the profile closure.  Define
\[
 B=X\setminus\Gamma.
\]
Then $X$ is compact metrizable and $B$ is closed.  Every coordinate in
$\mathcal F$ extends continuously to $X$; we retain its name.  A margin
$\psi_O$ also extends, through its expression in base coordinates.
The invariant family makes left translation extend to a homeomorphism
of $X$: it is the restriction of a coordinate permutation, with the
inverse permutation induced by the inverse group element.

Enumerate $\mathcal F=\{F_1,F_2,\ldots\}$ and choose positive numbers
$c_j$ satisfying
\[
 \sum_j c_j(1+\|dF_j\|_2)<\infty;
\]
for example, take $c_j=2^{-j}/(1+\|dF_j\|_2)$.  Equip $X$ with the
compatible metric
\begin{equation}\label{compact:metric}
 \rho(x,y)=\sum_j c_j|F_j(x)-F_j(y)|,
 \qquad \lambda(e)=\rho(e^-,e^+).
\end{equation}
Uniform convergence gives continuity of $\rho$, and positivity of all
weights makes it induce the product topology.  Minkowski's inequality
for finite partial sums, followed by monotone convergence, yields
\begin{equation}\label{compact:edge-lengths}
 \|\lambda\|_{\ell^2(E(G))}
 \le\sum_j c_j\|dF_j\|_2<\infty.
\end{equation}
In particular, for any fixed $\epsilon>0$ only finitely many edges have
$\rho$-length at least $\epsilon$.

\begin{lemma}\label{compact:boundary}
The boundary $B$ is a nondegenerate compact connected metric space.
\end{lemma}

\begin{proof}
A closed subset of $X$ contained in its isolated vertices is finite,
since it is compact and discrete.  Since $\Gamma$ is infinite, this
first implies $B\ne\varnothing$.  If $B$ were a singleton $\{x\}$,
then, for each $\epsilon>0$, the vertex set where
$|f(v)-f(x)|\ge\epsilon$ would be finite.  This would make $f$ tend to
a constant at infinity, contrary to Lemma~\ref{compact:potential}.
Thus $B$ has at least two points.

If $B$ were disconnected, write it as a disjoint union of two
nonempty compact sets.  Choose neighborhoods of those sets in $X$
whose closures have positive distance apart.  Their complement is a
closed subset of $X$ disjoint from $B$, hence a finite vertex set.
Each neighborhood contains infinitely many vertices.  Only finitely
many graph edges join the two neighborhoods, by
\eqref{compact:edge-lengths}.  Delete their endpoints together with
the finite complement.  The remaining graph has infinitely many
vertices on each side and no path between the sides.  A connected
locally finite graph minus finitely many vertices has only finitely
many components, so each side contains an infinite component.  This
contradicts one-endedness of $G$.
\end{proof}

\subsection{Paths near a boundary point}

We have a connected boundary, but connectedness alone does not control
paths near one of its points.  The extra coordinates in
\eqref{compact:localized} give exactly the following control of vertex
paths.

\begin{lemma}\label{compact:near-path}
For every $x\in B$ and every open neighborhood $W$ of $x$ in $X$, there
is an open neighborhood $V\subset W$ of $x$ such that any two vertices
in $V$ can be joined by a finite graph path whose vertices all lie in
$W$.

Equivalently, if $u_n,v_n\in\Gamma$ tend to the same $x\in B$, there
are joining graph paths $P_n$ with
$\sup\{\rho(w,x):w\text{ is a vertex of }P_n\}\to0$.
\end{lemma}

\begin{proof}
For a good box $O$ with its finite component list $U_1,\ldots,U_r$,
the identities
\begin{equation}\label{compact:component-identities}
 \sum_{j=1}^r h_{O,U_j}=\psi_O,
 \qquad h_{O,U_i}h_{O,U_j}=0\quad(i\ne j)
\end{equation}
hold on vertices and therefore, by continuity, on all of $X$.
Consequently, if $\psi_O(x)>0$, exactly one component, denoted
$U_O(x)$, satisfies
\begin{equation}\label{compact:selected-component}
 h_{O,U_O(x)}(x)=\psi_O(x)>0.
\end{equation}
All sufficiently near vertices belong to this component.  This is
where the finiteness assertion in Lemma~\ref{compact:good-boxes} is
used.

Shrink $W$ to a neighborhood defined by finitely many coordinate
conditions, each allowing a positive tolerance about its value at
$x$.  The base coordinates define a continuous projection from $X$
to $[0,1]^{\mathcal F_0}$.  Choose a good basis box $O'$ about the base
profile of $x$, small enough to have the following properties.
First, every tested base coordinate satisfies its prescribed tolerance
throughout $O'$.  Second, for every tested $h_{O,U}$ with
$\psi_O(x)>0$, we have $O'\subset O$ and the variation of $\psi_O$
from $\psi_O(x)$ on $O'$ is smaller than the prescribed tolerance.
Third, if a tested box has $\psi_O(x)=0$, its margin is smaller than
that tolerance throughout $O'$.  These are finitely many open
conditions on the base profile, so Lemma~\ref{compact:good-boxes}
supplies such an $O'$.

We have $\psi_{O'}(x)>0$.  Let $U'=U_{O'}(x)$ be the component selected
by \eqref{compact:selected-component}.  For any tested box $O$ of
positive margin at $x$, the inclusion $O'\subset O$ and connectedness
of $U'$ place $U'$ in one component of $O$.  That component is
$U_O(x)$.  Indeed, density of the vertices in $X$ supplies a vertex
near $x$ at which the two coordinates selecting $U'$ and $U_O(x)$ are
both positive.

Every tested localized coordinate therefore satisfies its
condition on $U'$.  If its box has positive margin and its
component is $U_O(x)$, it equals $\psi_O$, whose variation was
controlled.  If its component differs from $U_O(x)$, it is identically
zero on $U'$ and at $x$.  This includes the case of zero localized
value with positive box margin.  Finally, a tested box of zero margin
satisfies $0\le h_{O,U}\le\psi_O$, so the third condition on $O'$
controls its value.  Thus all vertices of $U'$ lie in $W$.

The open set
\[
 V=W\cap\{y\in X:
       h_{O',U'}(y)>\tfrac12\psi_{O'}(x)\}
\]
contains $x$, and every vertex in it belongs to $U'$.  Joining two
such vertices inside the graph component $U'$ proves the first
assertion.  Applying that assertion to successively smaller metric
balls about $x$, and choosing their radii slowly enough for both
endpoint sequences, proves the sequential formulation.  Conversely,
failure for one fixed $W$ would give two sequences converging to $x$
with no joining paths in $W$, contradicting that formulation.
\end{proof}
